\documentclass{amsart}
% For dealing with references we use the comment environment
\usepackage{verbatim}
\newenvironment{reference}{\comment}{\endcomment}
%\newenvironment{reference}{}{}
\newenvironment{slogan}{\comment}{\endcomment}
\newenvironment{history}{\comment}{\endcomment}

% For commutative diagrams we use Xy-pic
\usepackage[all]{xy}

% We use 2cell for 2-commutative diagrams.
\xyoption{2cell}
\UseAllTwocells

% We use multicol for the list of chapters between chapters
\usepackage{multicol}

% This is generally recommended for better output
\usepackage{lmodern}
\usepackage[T1]{fontenc}

% For cross-file-references

% Package for hypertext links:
\usepackage{hyperref}

% For any local file, say "hello.tex" you want to link to please

% Theorem environments.
%
\theoremstyle{plain}
\newtheorem{theorem}[subsection]{Theorem}
\newtheorem{proposition}[subsection]{Proposition}
\newtheorem{lemma}[subsection]{Lemma}

\theoremstyle{definition}
\newtheorem{definition}[subsection]{Definition}
\newtheorem{example}[subsection]{Example}
\newtheorem{exercise}[subsection]{Exercise}
\newtheorem{situation}[subsection]{Situation}

\theoremstyle{remark}
\newtheorem{remark}[subsection]{Remark}
\newtheorem{remarks}[subsection]{Remarks}

\numberwithin{equation}{subsection}

% Macros
%
\def\lim{\mathop{\mathrm{lim}}\nolimits}
\def\colim{\mathop{\mathrm{colim}}\nolimits}
\def\Spec{\mathop{\mathrm{Spec}}}
\def\Hom{\mathop{\mathrm{Hom}}\nolimits}
\def\Ext{\mathop{\mathrm{Ext}}\nolimits}
\def\SheafHom{\mathop{\mathcal{H}\!\mathit{om}}\nolimits}
\def\SheafExt{\mathop{\mathcal{E}\!\mathit{xt}}\nolimits}
\def\Sch{\mathit{Sch}}
\def\Mor{\mathop{\mathrm{Mor}}\nolimits}
\def\Ob{\mathop{\mathrm{Ob}}\nolimits}
\def\Sh{\mathop{\mathit{Sh}}\nolimits}
\def\NL{\mathop{N\!L}\nolimits}
\def\CH{\mathop{\mathrm{CH}}\nolimits}
\def\proetale{{pro\text{-}\acute{e}tale}}
\def\etale{{\acute{e}tale}}
\def\QCoh{\mathit{QCoh}}
\def\Ker{\mathop{\mathrm{Ker}}}
\def\Im{\mathop{\mathrm{Im}}}
\def\Coker{\mathop{\mathrm{Coker}}}
\def\Coim{\mathop{\mathrm{Coim}}}

% Boxtimes
%
\DeclareMathSymbol{\boxtimes}{\mathbin}{AMSa}{"02}

%
% Macros for moduli stacks/spaces
%
\def\QCohstack{\mathcal{QC}\!\mathit{oh}}
\def\Cohstack{\mathcal{C}\!\mathit{oh}}
\def\Spacesstack{\mathcal{S}\!\mathit{paces}}
\def\Quotfunctor{\mathrm{Quot}}
\def\Hilbfunctor{\mathrm{Hilb}}
\def\Curvesstack{\mathcal{C}\!\mathit{urves}}
\def\Polarizedstack{\mathcal{P}\!\mathit{olarized}}
\def\Complexesstack{\mathcal{C}\!\mathit{omplexes}}
% \Pic is the operator that assigns to X its picard group, usage \Pic(X)
% \Picardstack_{X/B} denotes the Picard stack of X over B
% \Picardfunctor_{X/B} denotes the Picard functor of X over B
\def\Pic{\mathop{\mathrm{Pic}}\nolimits}
\def\Picardstack{\mathcal{P}\!\mathit{ic}}
\def\Picardfunctor{\mathrm{Pic}}
\def\Deformationcategory{\mathcal{D}\!\mathit{ef}}

\setlength{\emergencystretch}{3em}
\begin{document}
\title[Stacks Project, Tag 07RT]{465 --- Pushouts of scheme thickenings along affine morphisms exist as schemes}
\maketitle
\noindent Copyright (C) 2005--2025 Johan de Jong. Stacks Project authors. Source: \href{https://stacks.math.columbia.edu/tag/07RT}{Tag 07RT}.

\noindent Editorial difficulty note (not author text). Difficulty H2: The fibre product of structure sheaves must first be proved to define a scheme, not merely a ringed-space pushout. The proof identifies its local affine rings and localized sections through a nilpotent fibre-product calculation, then glues the local universal mapping properties to establish the global pushout..

\noindent Editorial provenance note (not author text). History: excerpt from revision \texttt{a0444\allowbreak 6e57e\allowbreak c1fbc\allowbreak 252a8\allowbreak 71afc\allowbreak ec775\allowbreak 2fb28\allowbreak 07b14}, more-morphisms.tex, lines 4206--4269. Reference presentation was adapted; original-author context and core support blocks were retained or added. The mathematical wording is unchanged. Licensed under GNU FDL 1.2 or later; no invariant sections or cover texts. See ../licenses/COPYING. Context blocks reproduce author definitions and supporting results; they are not additional counted problems.

\begin{definition}
\label{definition-thickening}
Thickenings.
\begin{enumerate}
\item We say a scheme $X'$ is a {\it thickening} of a scheme $X$ if
$X$ is a closed subscheme of $X'$ and the underlying topological spaces
are equal.
\item We say a scheme $X'$ is a {\it first order thickening} of a scheme $X$ if
$X$ is a closed subscheme of $X'$ and the quasi-coherent sheaf of ideals
$\mathcal{I} \subset \mathcal{O}_{X'}$ defining $X$ has square zero.
\item We say a scheme $X'$ is a {\it finite order thickening} of a scheme $X$
if $X$ is a closed subscheme of $X'$ and the quasi-coherent sheaf of ideals
$\mathcal{I} \subset \mathcal{O}_{X'}$ defining $X$ is nilpotent, i.e.,
there exists an integer $n \geq 0$ such that $\mathcal{I}^{n + 1} = 0$.
\item We say a scheme $X'$ is an {\it $n$th order thickening} of a scheme $X$
if $X$ is a closed subscheme of $X$ and $\mathcal{I}^{n + 1} = 0$
where $\mathcal{I} \subset \mathcal{O}_{X'}$ is the
quasi-coherent sheaf of ideals defining $X$.
\item Given two thickenings $X \subset X'$ and $Y \subset Y'$ a
{\it morphism of thickenings} is a morphism $f' : X' \to Y'$ such that
$f'(X) \subset Y$, i.e., such that $f'|_X$ factors through the closed
subscheme $Y$. In this situation we set $f = f'|_X : X \to Y$ and we say
that $(f, f') : (X \subset X') \to (Y \subset Y')$ is a morphism of
thickenings.
\item Let $S$ be a scheme. We similarly define {\it thickenings over $S$}, and
{\it morphisms of thickenings over $S$}. This means that the schemes
$X, X', Y, Y'$ above are schemes over $S$, and that the morphisms
$X \to X'$, $Y \to Y'$ and $f' : X' \to Y'$ are morphisms over $S$.
\end{enumerate}
\end{definition}

\begin{lemma}
\label{lemma-pushout-along-thickening}
Let $X \to X'$ be a thickening of schemes and let $X \to Y$ be an affine
morphism of schemes. Then there exists a pushout
$$
\xymatrix{
X \ar[r] \ar[d]_f
&
X' \ar[d]^{f'}
\\
Y \ar[r]
&
Y'
}
$$
in the category of schemes. Moreover, $Y \subset Y'$ is a
thickening, $X = Y \times_{Y'} X'$, and
$$
\mathcal{O}_{Y'} = \mathcal{O}_Y \times_{f_*\mathcal{O}_X} f'_*\mathcal{O}_{X'}
$$
as sheaves on $|Y| = |Y'|$.
\end{lemma}

\begin{proof}
We first construct $Y'$ as a ringed space. Namely, as topological
space we take $Y' = Y$. Denote $f' : X' \to Y'$ the map of topological
spaces which equals $f$. As structure sheaf $\mathcal{O}_{Y'}$ we take
the right hand side of the equation of the lemma. To see that
$Y'$ is a scheme, we have to show that any point has an affine
neighbourhood. Since the formation of the fibre product of sheaves
commutes with restricting to opens, we may assume $Y$ is affine.
Then $X$ is affine (as $f$ is affine) and $X'$ is affine as well
(see Lemma \href{https://stacks.math.columbia.edu/tag/06AD}{lemma thickening affine scheme (Tag 06AD)}).
Say $Y \leftarrow X \rightarrow X'$ corresponds
to $B \rightarrow A \leftarrow A'$. Set $B' = B \times_A A'$; this
is the global sections of $\mathcal{O}_{Y'}$. As $A' \to A$ is surjective
with locally nilpotent kernel we see that $B' \to B$ is surjective
with locally nilpotent kernel. Hence $\Spec(B') = \Spec(B)$ (as
topological spaces). We claim that $Y' = \Spec(B')$. To see this
we will show for $g' \in B'$ with image $g \in B$ that
$\mathcal{O}_{Y'}(D(g)) = B'_{g'}$. Namely, by
More on Algebra, Lemma \href{https://stacks.math.columbia.edu/tag/01Z8}{lemma diagram localize (Tag 01Z8)} we see that
$$
(B')_{g'} = B_g \times_{A_h} A'_{h'}
$$
where $h \in A$, $h' \in A'$ are the images of $g'$. Since
$B_g$, resp.\ $A_h$, resp.\ $A'_{h'}$ is equal to $\mathcal{O}_Y(D(g))$,
resp.\ $f_*\mathcal{O}_X(D(g))$, resp.\ $f'_*\mathcal{O}_{X'}(D(g))$ the
claim follows.

\medskip\noindent
It remains to show that $Y'$ is the pushout.
The discussion above shows the scheme $Y'$
has an affine open covering $Y' = \bigcup W'_i$
such that the corresponding opens
$U'_i \subset X'$, $W_i \subset Y$, and
$U_i \subset X$ are affine open.
Moreover, if $A'_i$, $B_i$, $A_i$ are the rings corresponding to
$U'_i$, $W_i$, $U_i$, then
$W'_i$ corresponds to $B_i \times_{A_i} A'_i$.
Thus we can apply Lemmas \href{https://stacks.math.columbia.edu/tag/0ET0}{lemma basic example pushout (Tag 0ET0)} and
\href{https://stacks.math.columbia.edu/tag/0BMP}{lemma pushout fpqc local (Tag 0BMP)} to conclude our construction is a pushout
in the category of schemes.
\end{proof}
\end{document}
